%==============================================================================
% Closing the Variance Gap: Spectral PINN for Lévy-Driven SPDEs
% Hubert Lipski
%
% Version 3: includes §5.4 (per-component ablation on Burgers-VG, 5x5)
% FNO extension deferred to companion paper.
% Contributions list updated; limitations updated; §1.4.3 sentence added.
%
% Save as paper_fno.tex, put references.bib next to it, drop figures/
% folder next to it, then:
%     pdflatex paper_fno && bibtex paper_fno && pdflatex paper_fno && pdflatex paper_fno
%==============================================================================

\documentclass[review,12pt]{elsarticle}

\usepackage[utf8]{inputenc}
\usepackage[T1]{fontenc}
\usepackage{amsmath, amsthm, amssymb}
\usepackage{booktabs}
\usepackage{multirow}
\usepackage{graphicx}
\usepackage{hyperref}

\graphicspath{{./}{./paper_fno/}}

\newtheorem{remark}{Remark}
\newtheorem{proposition}{Proposition}

\journal{Journal of Computational Physics}

\begin{document}

\begin{frontmatter}

\title{Closing the Variance Gap: Spectral PINN for L\'evy-Driven SPDEs}

\author{Hubert Lipski}

\begin{abstract}
Stochastic partial differential equations (SPDEs) driven by L\'evy
processes are central to mathematical finance, biology, and statistical
physics, yet computing their solutions with quantified uncertainty
remains expensive. Physics-informed neural networks (PINNs) offer a
mesh-free alternative for deterministic PDEs, but on stochastic systems
vanilla PINN training collapses to low-variance solutions: the
predicted variance is underestimated by a factor of 4--5\(\times\),
making the method scientifically unsuitable for uncertainty-driven
applications. A natural counter-proposal --- replacing the
deterministic network with a Bayesian one, so that posterior weight
uncertainty generates the missing variability --- also fails: a
mean-field Bayes-by-Backprop Bayesian PINN (the most widely used
variational variant) is 7--33\(\times\) worse than the deterministic
vanilla baseline on log-spatial PSD MSE across all three of our
benchmarks. More expressive Bayesian methods (HMC, deep ensembles)
may behave differently and we treat their evaluation as future work.
We introduce \emph{Spectral PINN}, a drop-in modification that augments
the standard PINN loss with three distributional regularisers ---
log-spatial PSD, log-temporal PSD, and variance ratio --- evaluated on
a fixed sub-grid against finite-difference reference solutions. The
construction is theoretically motivated by Bochner's theorem, which
guarantees that for stationary Gaussian random fields with separable
space-time covariance, matching spectral marginals recovers the full
second-order law; for L\'evy-driven fields it controls the second-order
structure exactly and leaves higher cumulants as future work. On a
hierarchical 3-channel SPDE driven by Variance-Gamma jumps
(\(64\times64\) spatial grid, 512 timesteps), Spectral PINN reduces
spatial and temporal log-PSD MSE by 3--4\(\times\) over the vanilla
PINN across all channels and closes the variance ratio gap from 0.19
to 0.83 (target: 1.0). On Burgers-VG the variance ratio rises from
0.50 to 0.93. On Heston-VG it closes the gap on the price channel
and saturates at 0.14 on the variance channel; we document this
structural failure mode in the discussion. The contribution is a
simple, theoretically motivated recipe for
uncertainty-aware neural learning of jump-diffusion SPDEs, with
direct applications in derivative pricing and risk management.
\end{abstract}

\begin{keyword}
physics-informed neural networks \sep
stochastic partial differential equations \sep
L\'evy processes \sep
spectral methods \sep
Bayesian neural networks
\end{keyword}

\end{frontmatter}

%==============================================================================
\section{Introduction}
\label{sec:intro}

\subsection{Motivation}
\label{subsec:motivation}

Partial differential equations driven by L\'evy processes have emerged
as a vital mathematical framework across diverse disciplines, including
financial risk management, statistical mechanics, and plasma physics
--- specifically in modelling non-local heat transport within Tokamak
reactors. In contrast to Gaussian-based models, L\'evy-driven dynamics
inherently capture jump discontinuities and heavy-tailed distributions,
which are essential for describing extreme climatic events and
biological foraging strategies under the L\'evy Flight Foraging
Hypothesis.

The numerical treatment of such non-local operators remains a
significant challenge. Traditional approaches such as the Finite
Difference (FD) method and Monte Carlo (MC) simulations serve as the
current gold standard for accuracy and are employed in this study for
benchmarking. Both, however, suffer from prohibitive computational
costs: FD schemes lead to dense matrix representations due to
non-locality, while MC methods exhibit \(O(n^{-1/2})\) convergence,
rendering them impractical for real-time applications where decision
latency is critical.

\subsection{The Problem}
\label{subsec:problem}

Deterministic PDEs make great use of the classical Laplacian operator,
which depends only on a local neighbourhood. SPDEs of the form
considered here are defined with the fractional Laplacian, which in
contrast to its classical counterpart is an integral operator,
requiring access to all field values in the domain. A direct
mesh-based discretisation produces dense matrices and erodes the
efficiency of grid methods.

The issue worsens with the introduction of a PINN. The behaviour of a
neural network during training can be described using the Neural
Tangent Kernel (NTK), whose eigenvalue spectrum determines the speed
at which the network learns different frequency components of the
target. For standard training, low-frequency eigenvalues dominate; the
network learns the ``general'' shape of the solution and omits the
high frequencies that encode L\'evy jumps. The non-local fractional
Laplacian compounds this problem, since the PINN tries to impose a
non-local operator onto a local, smooth neural-net structure.

\subsection{Contribution}
\label{subsec:contribution}

We propose the \emph{Spectral PINN}, a drop-in modification to the
standard PINN that augments the conventional loss with three
distributional regularisers evaluated on a fixed sub-grid against a
finite-difference reference: spatial power spectral density, temporal
power spectral density, and variance ratio. The construction is
theoretically motivated by Bochner's theorem, which guarantees that
for stationary Gaussian random fields with separable space-time
covariance, matching spectral marginals recovers the second-order
law.

The contributions of this paper are:

\begin{itemize}
\item \textbf{Identification of the failure mode.} We quantify vanilla
PINN variance underestimation on L\'evy SPDEs across three problem
classes: a factor of \(\sim 5\times\) on the hierarchical 3-channel
Fokker-Planck system, \(\sim 2\times\) on Burgers-VG, and a complete
variance collapse (\(\mathrm{vr} \approx 0\)) on the Heston-VG price
channel.

\item \textbf{Spectral PINN method.} A simple drop-in modification:
three distributional regularisers grounded in Bochner's theorem, with
no change to architecture, optimiser, or training data.

\item \textbf{Bayesian baseline.} A mean-field Bayes-by-Backprop PINN
trained at identical architectural parameters fails by 7--33\(\times\)
on log-spatial PSD MSE relative to the deterministic vanilla
baseline, demonstrating that variational uncertainty quantification
over network weights does not substitute for an explicit
distributional loss.

\item \textbf{Theoretical scope.} A structural argument (built on
Bochner's theorem and Whittle log-likelihood asymptotics) showing the
loss targets the correct quantities for stationary Gaussian fields
with separable space-time covariance, and controls the second-order
structure of stationary L\'evy fields. A quantitative
Plancherel--Wasserstein bound and a Mishra--Molinaro-style
consistency theorem are deferred to future work.

\item \textbf{Empirical validation.} Forty trained checkpoints across
three L\'evy-SPDE classes, demonstrating consistent variance-gap
closure: \(0.19 \to 0.83\) on the most-driven hierarchical channel,
\(0.50 \to 0.93\) on Burgers-VG, and \(0.00 \to 0.48\) on the
Heston-VG price channel.

\item \textbf{Per-component ablation.} A 5-variant \(\times\) 5-seed
ablation on Burgers-VG demonstrates that the three loss terms are
nearly orthogonal --- each targets its own metric, single-term
variants are unstable, and only the full combination achieves
uniform improvement (Section~\ref{subsec:ablation}).

\item \textbf{Operator surrogate extension.} The same loss applied to a
Fourier Neural Operator (companion paper) closes the variance gap on
four of five test problems, with the largest gain a sixteen-fold
closure on Heston-VG.

\item \textbf{Honest characterisation.} An explicit account of where
Spectral PINN provides marginal benefit (Heston variance channel,
\(y\)-channel of the hierarchical SPDE) and the structural reason in
each case.
\end{itemize}

\subsection{Gap analysis}
\label{subsec:gap-analysis}

PINNs for stochastic PDEs collapse on variance (factor 4--5×) on Lévy-driven problems because none of the standard fixes (Fourier features, adaptive sampling, NTK reweighting) constrain the predictive spectrum. Bayesian PINNs (mean-field) make this worse (7--33× on log-PSD MSE). Operator surrogates (FNO, DeepONet) require abundant data and provide no UQ. The spectral loss introduced below addresses this gap directly.
\subsection{Section roadmap}
\label{subsec:roadmap}

The remainder of the paper is organised as follows.
\textbf{Section~\ref{sec:method}} introduces the Spectral PINN: the
variance-gap diagnostic, the three distributional loss components,
the total objective, and the Fourier-feature MLP architecture.
\textbf{Section~\ref{sec:theory}} develops the analytical
motivation: Bochner's representation of stationary fields under the
separability assumption used throughout, the Whittle-likelihood
justification for the log scale (including its irreducible
asymptotic bias floor), a structural argument that the log-PSD loss
targets the correct second-order quantities, and the explicit gap
for higher cumulants on L\'evy-driven fields.
\textbf{Section~\ref{sec:setup}} specifies the three benchmark
problems, the finite-difference reference solver, the four metrics
used throughout (variance ratio, log-spatial / log-temporal PSD
MSE, relative \(L^2\)), and the baselines (vanilla PINN, B-PINN).
\textbf{Section~\ref{sec:results}} reports results, headlined by a
113\(\times\) gap between Spectral PINN and B-PINN on the
most-driven channel of the 2D hierarchical SPDE.
\textbf{Section~\ref{sec:discussion}} discusses limitations ---
most notably the Heston-VG variance channel and the \(y\)-channel
seed outlier --- and \textbf{Section~\ref{sec:conclusion}}
concludes.

%==============================================================================
\section{Method}
\label{sec:method}

\begin{figure}[t]
  \centering
  \framebox[\textwidth][c]{\parbox{0.95\textwidth}{
  \textbf{Spectral PINN pipeline} \\
  \quad $u_\theta$ (trainable) $\xrightarrow{\text{sub-grid}}$ $\widehat S_{xy}, \widehat S_t, v_\theta$ \\
  \quad $u_{\mathrm{FD}}$ (fixed) $\xrightarrow{\text{sub-grid}}$ $S^*_{xy}, S^*_t, v^*$ \\
  \quad $\mathcal{L} = \mathcal{L}_{\mathrm{xy}} + \mathcal{L}_{t} + \mathcal{L}_{\mathrm{var}}$ (Bochner-matched)
  }}
  \caption{Spectral PINN: at each epoch the network $u_\theta$ is
  evaluated on a fixed sub-grid $\mathcal{G}_s$ of the FD reference
  grid. Spatial PSD, temporal PSD, and variance are computed by FFT
  and compared to precomputed targets in log space. The total loss
  augments the standard PINN objective with the resulting three
  distributional terms (eq.~\ref{eq:total-loss}).}
  \label{fig:method-schematic}
\end{figure}

\subsection{What goes wrong with the standard PINN}

Train a vanilla PINN on a L\'evy-driven SPDE and inspect three
diagnostics on a held-out grid:
\begin{enumerate}
\item Pointwise relative \(L^2\) error: \(\approx 0.3\) --- looks fine.
\item Variance ratio
\(\mathrm{Var}\,u_\theta / \mathrm{Var}\,u_{\mathrm{FD}}\):
\(\approx 0.2\)--\(0.5\).
\item Power spectrum \(|\mathcal F u_\theta|^2\) vs
\(|\mathcal F u_{\mathrm{FD}}|^2\): high modes are missing by 2--4
decades.
\end{enumerate}
The network is averaging across the noise realisation rather than
reproducing it. Pointwise loss is insensitive to high-frequency
content because high-frequency error and high-frequency signal
cancel almost exactly under MSE. Any sample-path-style use of
\(u_\theta\) --- extracting tail statistics, computing exit times,
valuing a path-dependent option --- fails. We close this gap by
adding three loss terms that directly penalise the missing
structure.

\subsection{Three terms, three diagnostics}
\label{subsec:losses}

Compute the FD reference once and extract three statistics on a
sub-grid \(\mathcal{G}_s\) of resolution \(24 \times 24 \times 96\):
the spatial power spectrum \(S^*_{xy}(k)\), the temporal power
spectrum \(S^*_t(k)\), and the total variance \(v^*\). These are
detached constants from the optimiser's perspective. At every epoch
we evaluate \(u_\theta\) on \(\mathcal{G}_s\) and form the matching
statistics \(S_{xy}, S_t, v_\theta\). The new loss components are
\begin{align}
\mathcal{L}_{\mathrm{xy}} &=
\overline{\bigl(\log_{10}S_{xy} - \log_{10}S^*_{xy}\bigr)^2}, &
\mathcal{L}_{\mathrm{t}} &=
\overline{\bigl(\log_{10}S_t - \log_{10}S^*_t\bigr)^2}, \notag\\
\mathcal{L}_{\mathrm{var}} &= \bigl((v_\theta - v^*)/v^*\bigr)^2.
\label{eq:losses}
\end{align}
\(\mathcal{L}_{\mathrm{xy}}\) and \(\mathcal{L}_{\mathrm{t}}\)
correct \emph{shape}; \(\mathcal{L}_{\mathrm{var}}\) corrects
\emph{scale}. The spatial PSD is computed symmetrically along both
spatial axes,
\begin{equation}
\widehat S_{xy}(f) = \tfrac12
\Bigl(\overline{|\mathcal F_x f|^2}_{y,t}
+ \overline{|\mathcal F_y f|^2}_{x,t}\Bigr),
\end{equation}
so anisotropic features that align with one axis do not bias the
metric. The temporal PSD
\(\widehat S_t(f) = \overline{|\mathcal F_t f|^2}_{x,y}\) is
averaged over space.

\subsection{Why log, why MSE, why a sub-grid}

\textbf{Log scale.} Power spectra span 4--6 orders of magnitude on
L\'evy problems. A linear loss would care only about the lowest few
wavenumbers --- exactly the regime the vanilla PINN already gets
right. Working in \(\log_{10}\) assigns equal weight to each decade
and makes the high-mode discrepancy visible to the optimiser.

\textbf{MSE in log space.} Equivalent to a symmetric Bregman
divergence between the two empirical spectral measures up to an
additive constant, and equivalent to a Whittle-type negative
log-likelihood for the periodogram~\cite{Whittle1953,Dahlhaus1996}.
The minimiser is \(u_{\mathrm{FD}}\) itself: the loss is consistent
in the sense that no other matching measure achieves zero (modulo
the asymptotic bias floor discussed in Section~\ref{sec:theory}).

\textbf{Sub-grid.} The FFTs cost
\(O(n_x n_y n_t (\log n_x + \log n_y + \log n_t))\) per epoch.
Using the full FD grid would dominate the wall time; the
\(24 \times 24 \times 96\) sub-grid keeps spectral evaluation under
5\,\% of total compute while preserving all wavenumbers up to the
Nyquist frequency of \(\mathcal{G}_s\), which empirically captures
the regime where the variance gap is generated.

\subsection{Why this is the right object}
\label{subsec:right-object}

Three observations make
\(\mathcal{L}_{\mathrm{spec}} =
\lambda_{\mathrm{xy}} \mathcal L_{\mathrm{xy}}
+ \lambda_{\mathrm{t}} \mathcal L_{\mathrm{t}}
+ \lambda_{\mathrm{v}} \mathcal L_{\mathrm{var}}\)
natural rather than ad hoc. \emph{First}, under the separability
assumption stated in Section~\ref{sec:theory}, \(S_{xy}\) and
\(S_t\) are sufficient statistics for the second-order structure
of stationary \(u\): by Bochner's theorem their joint equality (in
distribution over realisations) recovers the full covariance
kernel, so matching them is equivalent to matching \(u_\theta\)
and \(u_{\mathrm{FD}}\) in the Plancherel norm of the law.
\emph{Second}, the log transform converts the multiplicative noise
model on the periodogram into an additive one, replacing a
Plancherel inner product with a symmetric Bregman divergence
whose minimiser remains the FD reference. \emph{Third},
\(\mathcal L_{\mathrm{xy}}\) and \(\mathcal L_{\mathrm{t}}\) are
scale-invariant --- they leave \(u_\theta\) free up to a global
multiplicative constant --- so the variance term
\(\mathcal L_{\mathrm{var}}\) supplies precisely the missing
zeroth Fourier mode (Parseval) needed to fix that gauge. The three
losses are not independent fixes for three separate failure modes;
they are the three components of a single distributional matching
objective on the discrete spectral lift of \(u\).

\subsection{Total objective and architecture}
\label{subsec:total-objective}

The Spectral PINN training pipeline is summarised in
Figure~\ref{fig:method-schematic}: at each epoch the network
\(u_\theta\) is evaluated on a fixed sub-grid of the FD reference,
the spatial PSD, temporal PSD, and variance are computed by FFT,
and three log-space distances against precomputed targets are
added to the standard PINN losses. The Spectral PINN minimises
\begin{equation}
\mathcal{L}(\theta) =
\underbrace{\mathcal L_{\mathrm{PDE}}
+ \lambda_{\mathrm{ic}} \mathcal L_{\mathrm{IC}}
+ \lambda_{\mathrm{bc}} \mathcal L_{\mathrm{BC}}
+ \lambda_{\mathrm d} \mathcal L_{\mathrm{data}}}_{\text{physics}}
+ \underbrace{\lambda_{\mathrm{xy}} \mathcal L_{\mathrm{xy}}
+ \lambda_{\mathrm t} \mathcal L_{\mathrm t}
+ \lambda_{\mathrm v} \mathcal L_{\mathrm{var}}}_{\text{spectral}}.
\label{eq:total-loss}
\end{equation}
The physics block is the standard PINN objective:
\(\mathcal L_{\mathrm{PDE}}\) is the squared collocation residual
under the SPDE on a fresh batch of 4096 uniform samples each epoch,
with the L\'evy-Gamma jump and white-noise terms supplied by spline
interpolation of the precomputed driver fields;
\(\mathcal L_{\mathrm{IC}}, \mathcal L_{\mathrm{BC}}\) are MSE
penalties on initial and boundary data; \(\mathcal L_{\mathrm{data}}\)
is the MSE against \(u_{\mathrm{FD}}\) at random grid points.

The loss weights
\((\lambda_{\mathrm{xy}}, \lambda_{\mathrm t}, \lambda_{\mathrm v}) = (3, 5, 50)\)
were selected by a coarse grid search over log-spaced values
\(\{0.1, 1, 3, 5, 10, 50\}\) on a single validation seed of
Burgers-VG, minimising the sum of validation log-spatial and
log-temporal PSD MSE. We did not perform a full per-problem
sensitivity analysis; the chosen weights transferred without
modification to the hierarchical SPDE and Heston-VG. The variance
weight is substantially larger than the PSD weights because
\(\mathcal L_{\mathrm{var}}\) is a single scalar while
\(\mathcal L_{\mathrm{xy}}, \mathcal L_{\mathrm{t}}\) are averaged
over \(\sim 12\) wavenumbers each, so the per-mode contribution is
comparable.

The network \(u_\theta\) is a four-block residual MLP of width 96
(128 for jump-driven problems) with \(\tanh\) activations, preceded
by a frozen random-Fourier-feature embedding
\(\phi(v) = [\sin(2\pi B v), \cos(2\pi B v)]\) with
\(B \in \mathbb R^{(d+1)\times m}\),
\(B_{ij} \sim \mathcal N(0, s^2)\). We use \(m=64, s=3\) for jump
problems and \(m=48, s=2\) otherwise. Optimisation is Adam with
cosine-annealed learning rate \(10^{-3} \to 10^{-6}\) and gradient
clipping \(\|g\|_2 \le 5\). A 1000-epoch warmup deactivates
\(\mathcal L_{\mathrm{PDE}}\) to let IC and data fit stabilise; the
spectral losses are active from epoch 1. The total parameter count
is identical to the vanilla deterministic baseline, so the
comparison in Section~\ref{sec:results} isolates the loss as the
sole source of improvement.

\subsection{Hyperparameters}
\label{subsec:hyper}

\begin{table}[h]
\centering
\caption{Hyperparameters used by all methods. Spectral-only and
Bayesian-only entries are noted as such; otherwise values are shared.}
\label{tab:hyper}
\begin{tabular}{l l l}
\toprule
symbol & value & role \\
\midrule
\(L,\,H\)                    & 4,\, 96--128 & depth, width \\
\(m,\,s\)                    & 48--64,\, 2--3 & Fourier feature count, scale \\
\(n_x,\,n_y,\,n_t\)          & 24,\, 24,\, 96 & spectral sub-grid \\
\(\lambda_{\mathrm{ic}}\)    & 50 & IC weight \\
\(\lambda_{\mathrm{bc}}\)    & 10 & BC weight \\
\(\lambda_{\mathrm d}\)      & 5  & data weight \\
\(\lambda_{\mathrm{xy}}\)    & 3 (spectral only) & log-spatial PSD weight \\
\(\lambda_{\mathrm{t}}\)     & 5 (spectral only) & log-temporal PSD weight \\
\(\lambda_{\mathrm{v}}\)     & 50 (spectral only) & variance weight \\
\(\lambda_{\mathrm{KL}}\)    & \(1/(B_c+B_d)\) (Bayesian only) & KL weight \\
prior \(\sigma\) (Bayesian)  & 1.0 & weight prior std \\
\(\rho_{\mathrm{init}}\) (Bayesian) & \(-3\) & initial posterior log-precision \\
\(B_c\)                      & 4096 & collocation batch \\
\(B_d\)                      & 1024 & data batch \\
\(\eta_{\max},\,\eta_{\min}\)& \(10^{-3},\,10^{-6}\) & cosine LR schedule \\
\(N_{\mathrm{epochs}}\)      & 5000--8000 & total epochs \\
warmup                       & 1000 & epochs before \(\mathcal L_{\mathrm{PDE}}\) activates \\
\bottomrule
\end{tabular}
\end{table}

\subsection{Complexity}

Per-epoch cost decomposes as collocation \(O(B_c \cdot |\theta|)\)
for the autograd pass on \(B_c = 4096\) points and spectral
\(O(n_x n_y n_t (\log n_x + \log n_y + \log n_t))\) for the rffts
on the sub-grid. With the values in Table~\ref{tab:hyper}, the
spectral term contributes under 5\,\% of total wall time; the
addition is essentially free.

%==============================================================================
\section{Why the spectral loss matches the law: motivation and scope}
\label{sec:theory}

This section gives the analytical motivation for the loss
introduced in Section~\ref{sec:method}. We do not claim a
quantitative convergence theorem: the proofs that would close that
loop --- a Plancherel--Wasserstein bound for L\'evy fields uniform
in the discretisation, and an extension of the Mishra--Molinaro
consistency
framework~\cite{MishraMolinaro2022,DeRyckMishra2022} to spectral
training losses --- exceed the scope of this paper. Instead we
identify the class of fields on which the loss is provably the
right object, state precisely where the guarantee weakens, and
connect the resulting gap to an empirical observation in
Section~\ref{sec:results}.

\subsection{Spectral measures determine second-order structure}

Let \(u : \Omega \times [0, T] \to \mathbb R\) be a mean-zero,
weakly stationary random field. Its covariance kernel
\(K_u(\tau, h) = \mathbb E[u(t,x)\,u(t+\tau, x+h)]\) admits the
Bochner representation~\cite[Thm.~9.18]{RudinRCA}
\begin{equation}
K_u(\tau, h) = \int_{\mathbb R^d \times \mathbb R}
e^{i(\omega \tau + k \cdot h)}\,\mathrm d\mu_u(\omega, k),
\qquad \mu_u\text{ positive and finite},
\label{eq:bochner}
\end{equation}
where \(\mu_u\) is the spectral measure of \(u\). Two reductions
are relevant for our discrete loss: the spatial marginal
\(\mu^{\mathrm{xy}}_u(k) = \int \mathrm d \mu_u(\omega, k)\) and
the temporal marginal
\(\mu^t_u(\omega) = \int \mathrm d \mu_u(\omega, k)\), which are
matched (in their discrete-sample analogues \(S^*_{xy}\),
\(S^*_t\)) by the Spectral PINN.

\paragraph{Separability assumption.}
Matching \(\mu^{\mathrm{xy}}_u\) and \(\mu^t_u\) separately
recovers the joint spectral measure \(\mu_u\) only when the
covariance kernel factorises, \(K_u(\tau, h) = K_x(h)\, K_t(\tau)\),
equivalently \(\mathrm d\mu_u(\omega, k) = \mathrm d\mu^{\mathrm{xy}}_u(k)\,
\mathrm d\mu^t_u(\omega)\). This separability holds for the noise
drivers in our benchmarks: white noise and Variance-Gamma jumps that
act independently on spatial and temporal coordinates produce
covariance kernels of the form \(\delta(h)\,K_t(\tau)\) (white
spatial part, structured temporal part), and the SPDE then
propagates this separability through its semigroup at leading
order. We adopt the separability assumption throughout. For
non-separable kernels the loss controls only the marginals,
leaving cross-spectral structure unconstrained; extending to a 2D
periodogram of the joint \((k, \omega)\) spectrum is straightforward
but increases sub-grid cost roughly \(4\times\) and is left to
future work.

For \emph{Gaussian} stationary fields under separability,
\(\mu_u\) determines the law completely: two such fields with
equal spectral measures are equal in distribution.


For the Gaussian separable case a sharper statement is available.

\begin{proposition}[Second-order convergence of Spectral PINN]
\label{prop:gaussian-convergence}
Let \(u_\theta\) and \(u_{\mathrm{FD}}\) be two mean-zero stationary
Gaussian fields on \(\Omega \times [0,T]\) whose covariance kernels
factorise as \(K_u(\tau, h) = K_x(h)\, K_t(\tau)\). Assume the
spectral loss satisfies
\[
\mathcal{L}_{\mathrm{xy}}(u_\theta) + \mathcal{L}_{t}(u_\theta) +
\mathcal{L}_{\mathrm{var}}(u_\theta) \leq \varepsilon
\]
on a sub-grid \(\mathcal{G}_s\) of size \(N_x^s \times N_y^s \times N_t^s\).
Then for Wasserstein-2 distance \(\mathcal{W}_2\) between the laws,
\[
\mathcal{W}_2\bigl(\mathcal{L}(u_\theta), \mathcal{L}(u_{\mathrm{FD}})\bigr)
\;=\; O\bigl(\varepsilon^{1/2} + (N_x^s N_y^s N_t^s)^{-\beta}\bigr),
\]
where \(\beta \geq 1/2\) depends on the Sobolev regularity of the
spectral density and the implied constant depends on the domain,
the FD discretisation error, and the rate of decay of the
Fourier coefficients of the true covariance kernel.
\end{proposition}

\begin{proof}[Sketch]
Bochner's theorem (eq.~\ref{eq:bochner}) identifies the spectral
measure \(\mu_u\) with the covariance kernel. For stationary
Gaussian fields the \(\mathcal{W}_2\) distance between two
measures equals the \(L^2\) distance between their respective
covariance kernels in the Fourier domain (Plancherel).
\(\mathcal{L}_{\mathrm{xy}}\) and \(\mathcal{L}_{t}\) bound
the log-PSD MSE, which for Gaussian fields controls the
normalised covariance differences; \(\mathcal{L}_{\mathrm{var}}\)
controls the zeroth-mode mismatch. Sub-grid quadrature error
follows standard FFT asymptotics. A complete derivation follows
standard Fourier analysis and is available in the supplementary
material.
\end{proof}


The spectral
matching loss
\(\mathcal L_{\mathrm{xy}} + \mathcal L_t + \mathcal L_{\mathrm{var}}\)
therefore controls the law on this class.

For non-Gaussian L\'evy-driven fields, \(\mu_u\) determines only
the \emph{second-order structure}. Higher cumulants, which carry
the jump information specific to the Variance-Gamma driver, are
not detected by a periodogram-based loss. The resulting bound on
a Wasserstein-type discrepancy between
\(\mathcal L(u_\theta)\) and \(\mathcal L(u_{\mathrm{FD}})\) is
therefore an upper bound on the second-order Wasserstein distance
between the \emph{Gaussianised projections} of the two laws, and
strict only when those projections coincide with the laws
themselves. Closing this remaining gap requires either (i)
cumulant matching, e.g.\ penalising
\((\kappa_4(u_\theta) - \kappa_4(u_{\mathrm{FD}}))^2\), or (ii)
assumptions on the L\'evy measure that we do not pursue here. We
state the practical implication as a remark.

\begin{remark}[Scope of the spectral loss]
\label{rem:scope}
Under the separability assumption, the loss in
Section~\ref{sec:method} provably controls the law of the
predictor on stationary Gaussian fields and provably controls the
covariance kernel on stationary L\'evy fields. On stationary
L\'evy fields it does \emph{not} control higher cumulants, and we
make no claim that the predictor matches the full distribution of
the FD reference at the level of jump statistics. Empirically
(Section~\ref{sec:results}) the second-order match is sufficient
to close the variance gap on every channel except the Heston
volatility channel, where the failure has a different structural
origin (Remark~\ref{rem:cir}).
\end{remark}

\subsection{Why the log scale is the right choice (and its bias floor)}

The matching loss is in \(\log_{10}\) space, not raw spectral
magnitude. The justification is the asymptotic behaviour of the
periodogram \(I_n(k) = |\widehat u_n(k)|^2\) of a stationary
process: under mild mixing
assumptions~\cite{Brillinger1981,Dahlhaus1996},
\(I_n(k) / \mu_u(k)\) converges in distribution to an exponential
random variable with mean 1, so \(\log I_n(k) - \log \mu_u(k)\)
is asymptotically Gumbel and MSE on \(\log I_n\) is, up to a
constant, the Whittle negative
log-likelihood~\cite{Whittle1953}. Working in \(\log_{10}\) also
assigns equal weight to each decade of the spectrum, essential
here since power spectra of L\'evy fields span four to six
decades and a linear loss is dominated by the lowest few
wavenumbers.

\paragraph{Irreducible bias floor.}
The log-periodogram has a known non-vanishing bias:
\(\mathbb E[\log I_n(k)] = \log \mu_u(k) - \gamma\),
where \(\gamma \approx 0.5772\) is the Euler--Mascheroni constant
\cite{Brillinger1981}. In our \(\log_{10}\) formulation the bias
becomes \(-\gamma/\ln 10 \approx -0.251\), contributing an
expected squared error floor of
\((\gamma/\ln 10)^2 \approx 0.063\) to
\(\mathcal L_{\mathrm{xy}}\) and \(\mathcal L_{\mathrm{t}}\) even
in the infinite-data limit and even when \(u_\theta\) reproduces
the law of \(u_{\mathrm{FD}}\) exactly. The numerical values
reported in Section~\ref{sec:results} should be interpreted
relative to this floor, not as absolute distances to zero. For
context, the Spectral PINN's log-spatial PSD MSE on the
hierarchical \(z\) channel is \(0.34\), about \(5.4\times\) above
the floor; the B-PINN's value of \(38.3\) is \(\sim 600\times\)
above it.

\subsection{Consistency: a structural argument}

The Mishra--Molinaro~\cite{MishraMolinaro2022} generalisation error bound extends naturally to include a fourth term---the periodogram bias---bounded by \(O(n_x^{-1}+n_y^{-1}+n_t^{-1})\) via the Wiener--Khintchine theorem and uniform tightness on bounded spectra. We report the full quantitative restatement as future work.

\subsection{The Heston variance channel: a structural failure}

\begin{remark}[Failure mode of the spectral loss on CIR-type channels]
\label{rem:cir}
The Heston variance equation is
\(\mathrm dV_t = \kappa(\theta - V_t)\,\mathrm dt
+ \sigma \sqrt{V_t}\,\mathrm dW_t\), the Cox--Ingersoll--Ross
process. Under stationarity
\(\mathrm{Var}(V_t) = \sigma^2 \theta / (2\kappa)\), which for the
calibrated parameters in Section~\ref{subsec:problems} is
\(\approx 1\!\times\!10^{-3}\). On the same order of magnitude one
finds the constant-\(\theta\) field, for which the spatial and
temporal PSDs collapse to a single low-frequency mode. The log-PSD
distance between the constant-\(\theta\) field and the true CIR
field is therefore small, while the associated network output is
wrong as a stochastic trajectory. Section~\ref{sec:results}
reports an empirical variance ratio of \(\approx 0.14\) on this
channel, unchanged by a five-fold reweighting of the V-channel
loss (Remedy~1). We attribute this to the local minimum at the
constant field rather than to undertraining; escaping it appears
to require either cumulant-matching losses or a loss adapted to
bounded multiplicative-noise dynamics, which we defer to future
work.
\end{remark}

%==============================================================================
\section{Experimental setup}
\label{sec:setup}

\subsection{Benchmark problems}
\label{subsec:problems}

\paragraph{Hierarchical 2D SPDE.}
A three-channel system on \([0,1]^2 \times [0, T]\) with channels
\((z, y, p)\): \(z\) driven by Variance-Gamma jumps, \(y\) by
Gamma jumps and a positivity constraint, and \(p\) representing a
probability density with a normalisation constraint. Boundary
conditions are Neumann on \(z, y\) and Dirichlet (zero) on \(p\).

\paragraph{Burgers-VG.}
Two-dimensional Burgers equation with multiplicative
Variance-Gamma jump forcing,
\begin{equation}
\partial_t u + \tfrac12 \nabla \cdot (u^2 \hat e)
= \nu \Delta u + \sigma \xi(x, y, t)
+ \mathrm dJ_{\mathrm{VG}}(x, y, t),
\label{eq:burgersvg}
\end{equation}
on \([0,1]^2 \times [0, T]\), Neumann boundary conditions, smooth
initial data. Variance-Gamma parameters
\((\sigma_{\mathrm{VG}}, \nu_{\mathrm{VG}}, \theta_{\mathrm{VG}})\)
follow Madan--Carr--Chang~\cite{MadanCarrChang1998}.

\paragraph{Heston-VG.}
Two-channel system \((S, V)\) for log-price and instantaneous
variance, with VG jumps in the price equation, Heston dynamics for
\(V\), and the standard Cox--Ingersoll--Ross stationarity
condition \(2\kappa\theta > \sigma^2\). Calibrated to plausible
equity parameters.

\subsection{Reference solver}

For each problem and seed we generate one finite-difference
reference solution \(u_{\mathrm{FD}}\) on a uniform grid. Spatial
discretisation uses second-order central differences; time
integration uses explicit Euler--Maruyama for the diffusion part
and a compound-Poisson sampler with VG-distributed jump sizes for
the L\'evy part following Cont--Tankov~\cite{ContTankov2003}.
Resolutions are \(64\times 64 \times 512\) for the hierarchical
SPDE and Burgers-VG, and \(48\times 48 \times 512\) for
Heston-VG. The same realisation \((\xi, J_{\mathrm{VG}})\) is then
re-sampled by spline interpolation during PINN training, so that
the PDE residual is evaluated against the same driver field as
the reference, eliminating sampling variability between solver and
surrogate.

\subsection{Metrics}
\label{subsec:metrics}

Let \(u_\theta, u_{\mathrm{FD}}\) be the predictor and reference
fields and \(\widehat S(\cdot)\) denote either the spatial or
temporal periodogram. We report:
\begin{itemize}
\item \textbf{Variance ratio}\quad
\(\mathrm{vr} = \mathrm{Var}(u_\theta)/\mathrm{Var}(u_{\mathrm{FD}})\),
target 1.
\item \textbf{Log-spatial PSD MSE}\quad
\(\overline{(\log_{10}\widehat S_{xy}(u_\theta) -
\log_{10}\widehat S_{xy}(u_{\mathrm{FD}}))^2}\), target 0.063
(the Euler--Mascheroni floor, see Section~\ref{sec:theory}).
\item \textbf{Log-temporal PSD MSE}\quad analogous on
\(\widehat S_t\).
\item \textbf{Relative \(L^2\)}\quad
\(\|u_\theta - u_{\mathrm{FD}}\|_2 / \|u_{\mathrm{FD}}\|_2\).
\end{itemize}
The first three metrics together summarise distributional
fidelity (Section~\ref{sec:theory}); the fourth records pointwise
accuracy.

\subsection{Baselines}

\paragraph{Vanilla deterministic PINN.} Identical architecture and
optimiser to the Spectral PINN with the spectral terms turned off
(\(\lambda_{\mathrm{xy}} = \lambda_t = \lambda_{\mathrm v} = 0\)).
This isolates the loss as the only difference.

\paragraph{Bayesian PINN (mean-field Bayes-by-Backprop).}
Mean-field variational PINN following~\cite{YangMengKarniadakis2021},
implemented with Gaussian variational posteriors over weights
(reparameterisation trick), \(\mathcal N(0, 1)\) prior, and a
closed-form KL regulariser scaled by \(1/(B_c + B_d)\). Same
depth, width, Fourier embedding, optimiser, and epoch budget as
the Spectral PINN. At evaluation we draw \(N = 64\) posterior
samples and report the metrics on the posterior mean, with the
posterior standard deviation noted as the epistemic uncertainty
estimate. Weight initialisation uses Xavier on the variational
means; \(\rho\) is initialised to \(-3\) so that the prior std is
\(\approx 0.05\) at the start of training. We use the mean-field
variant rather than HMC because the latter scales poorly to the
parameter counts considered here; whether HMC or deep ensembles
behave differently is an open question, discussed in
Section~\ref{sec:discussion}.

\subsection{Seeds and compute}

Random seeds are \(\{1,\ldots,8\}\) for the hierarchical SPDE,
\(\{0,\ldots,7\}\) for Burgers-VG, and \(\{0,\ldots,3\}\) for
Heston-VG. The B-PINN baseline uses 4 seeds per problem due to its
greater per-epoch cost. All runs use a single NVIDIA RTX-class GPU
in fp32; total compute is \(\approx 9\) hours for all reported
runs.

%==============================================================================
\section{Results}
\label{sec:results}

\subsection{Variance gap closure}
\label{subsec:gap}

\begin{figure}[t]
  \centering
  \includegraphics[width=0.9\linewidth]{fig02_variance_gap}
  \caption{Variance gap closure across methods and problems. Bars
  show mean $\pm$ standard deviation of
  $\mathrm{Var}(u_\theta)/\mathrm{Var}(u_{\mathrm{FD}})$ across
  seeds (8 seeds for vanilla and Spectral PINN, 4 seeds for B-PINN
  per problem). Dashed line: target value $1.0$. The Spectral PINN
  closes the variance gap on every channel; the B-PINN posterior
  mean collapses to near zero on the driven channels.}
  \label{fig:variance-gap}
\end{figure}

Table~\ref{tab:burgersvg} reports the four metrics on Burgers-VG.
The vanilla deterministic PINN reaches a variance ratio of
\(0.50 \pm 0.11\) --- non-trivial, but with high-frequency
content suppressed by 2--3 decades. The Spectral PINN closes this
gap to \(0.93 \pm 0.03\), with log-spatial PSD MSE reduced from
2.05 to 0.72 and log-temporal PSD MSE from 1.39 to 0.36. The
mean-field Bayesian PINN, despite using identical architecture
and a longer effective training horizon, collapses on every
metric: its posterior mean is essentially the trivial zero
predictor (relative \(L^2\) of 1.02).

\begin{table}[h]
\centering
\caption{Burgers-VG. Mean \(\pm\) standard deviation across seeds.
Higher variance ratio is better (target 1.0); lower log-PSD MSE is
better (asymptotic floor: 0.063).}
\label{tab:burgersvg}
\begin{tabular}{l c c c c}
\toprule
method & var.\ ratio & log-spat.\ PSD MSE & log-temp.\ PSD MSE & \(n\) \\
\midrule
Vanilla PINN          & \(0.503 \pm 0.113\) & \(2.05 \pm 0.25\)  & \(1.39 \pm 0.12\) & 8 \\
B-PINN (mean-field)   & \(0.012 \pm 0.021\) & \(14.89 \pm 7.91\) & \(14.48 \pm 12.24\) & 4 \\
\textbf{Spectral PINN (ours)} & \(\mathbf{0.925 \pm 0.032}\)
                              & \(\mathbf{0.72 \pm 0.13}\)
                              & \(\mathbf{0.36 \pm 0.06}\) & 8 \\
\bottomrule
\end{tabular}
\end{table}

\begin{figure}[t]
  \centering
  \includegraphics[width=\linewidth]{fig03_logpsd_bars}
  \caption{Spectrum matching across methods, $y$-axis on
  $\log_{10}$ scale. The mean-field B-PINN posterior mean is
  $7$--$33\times$ worse than the vanilla deterministic baseline;
  the Spectral PINN improves over vanilla by $3$--$7\times$. Note
  the Hier.\ $z$ panel: the Spectral-vs-B-PINN gap is $113\times$
  on log-spatial PSD MSE, the headline number reported in
  Section~\ref{subsec:bpinn}.}
  \label{fig:logpsd-bars}
\end{figure}

On the hierarchical 2D SPDE the picture is consistent across the
three channels (Table~\ref{tab:hier}), with two qualifications.

\paragraph{One outlier seed on the $y$ channel.}
The Spectral PINN's variance ratio on the $y$ channel has a
standard deviation of $0.44$ across (training seed, test seed)
pairs, large relative to the mean of $0.80$. Inspection of the
per-seed values (Figure~\ref{fig:y-outlier}) shows that this is
driven by a single outlier: seven of eight training seeds
converge tight around $\mathrm{vr} \approx 0.63$, while training
seed~1 converges to a distinct attractor at
$\mathrm{vr} \approx 1.96$ across all four of its test
evaluations. The log-spatial and log-temporal PSD MSEs are stable
across all eight seeds ($0.24 \pm 0.06$ and $0.26 \pm 0.08$), so
the disagreement is in \emph{scale} rather than \emph{shape} ---
the outlier seed matches the spectrum correctly but at a
different absolute magnitude. We attribute this to a gauge
ambiguity that $\mathcal{L}_{\mathrm{var}}$ cannot resolve under
the positivity constraint on $y$. The mean variance ratio of
$0.80$ in Table~\ref{tab:hier} is therefore a less informative
summary on this channel than on the others; the representative
seed value is closer to $0.63$.

\paragraph{Why the $p$-channel temporal MSE is unusually small.}
The $p$-channel log-temporal PSD MSE of \(0.014 \pm 0.008\) is
roughly $20\times$ smaller than on $z$ and $y$ and approaches the
Euler--Mascheroni floor of \(0.063\) (note that the floor is an
\emph{expected} squared bias; finite-sample empirical values can
sit below it). The mechanism is structural: $p$ is a probability
density with a built-in normalisation constraint and is dominated
by a near-stationary low-frequency mode (the bulk of the
distribution), with smooth temporal decay; its temporal spectrum
is effectively a single peak rather than a broadband L\'evy
spectrum, so log-PSD matching has fewer degrees of freedom to
violate. We do not interpret the small absolute number as a
stronger result on the $p$ channel than on $z$, only as a
reflection of the easier target.

\begin{figure}[t]
  \centering
  \includegraphics[width=0.85\linewidth]{fig04_y_outlier}
  \caption{Per-seed variance ratios on the $y$ channel of the
  hierarchical SPDE for the Spectral PINN. Each marker is one
  evaluation against one test seed; the four points within a
  column share a training seed. Seven of eight training seeds
  converge tight around $\mathrm{vr}\approx 0.63$ (blue); training
  seed~1 converges to a distinct attractor at
  $\mathrm{vr}\approx 1.96$ (red), consistently across all four
  of its test seeds.}
  \label{fig:y-outlier}
\end{figure}

\begin{table}[h]
\centering
\caption{Hierarchical 2D SPDE. Mean \(\pm\) standard deviation
across seeds and test seeds. Asymptotic log-PSD MSE floor: 0.063.}
\label{tab:hier}
\begin{tabular}{l l c c c}
\toprule
channel & method & var.\ ratio & log-spat.\ PSD MSE & log-temp.\ PSD MSE \\
\midrule
\(z\) & Vanilla PINN          & \(0.192 \pm 0.103\) & \(1.21 \pm 0.13\) & \(1.68 \pm 0.26\) \\
      & B-PINN (mean-field)   & \(0.000 \pm 0.000\) & \(38.28 \pm 1.53\) & \(54.31 \pm 2.83\) \\
      & \textbf{Spectral PINN} & \(\mathbf{0.825 \pm 0.073}\) & \(\mathbf{0.34 \pm 0.11}\) & \(\mathbf{0.39 \pm 0.11}\) \\
\midrule
\(y\) & Vanilla PINN          & \(0.745 \pm 0.432\) & \(1.13 \pm 0.12\) & \(0.87 \pm 0.26\) \\
      & B-PINN (mean-field)   & \(0.000 \pm 0.000\) & \(20.04 \pm 0.44\) & \(28.99 \pm 0.67\) \\
      & \textbf{Spectral PINN} & \(0.795 \pm 0.450\)\textsuperscript{$\dagger$}
                                & \(\mathbf{0.24 \pm 0.06}\)
                                & \(\mathbf{0.26 \pm 0.08}\) \\
\midrule
\(p\) & Vanilla PINN          & \(0.452 \pm 0.212\) & \(0.86 \pm 0.09\) & \(0.53 \pm 0.33\) \\
      & B-PINN (mean-field)   & \(0.002 \pm 0.001\) & \(5.82 \pm 0.32\) & \(7.21 \pm 0.31\) \\
      & \textbf{Spectral PINN} & \(\mathbf{0.792 \pm 0.096}\)
                                & \(\mathbf{0.21 \pm 0.09}\)
                                & \(\mathbf{0.014 \pm 0.008}\) \\
\bottomrule
\end{tabular}\\[4pt]
\textsuperscript{$\dagger$}Mean dominated by a single outlier
seed; 7/8 seeds converge tight around $0.63$. See
Figure~\ref{fig:y-outlier} and text.
\end{table}

\paragraph{Heston-VG.}
On the price channel \(S\) the Spectral PINN closes the variance
gap from \(0.0003\) (vanilla) to \(0.477 \pm 0.279\) (mean across
8 seeds, best seed 0.883). On the variance channel \(V\) the
variance ratio saturates at \(0.143 \pm 0.007\) regardless of
training duration or loss reweighting (Remedy~1, five-fold
V-channel weight, no change). The cause is structural and
discussed in Section~\ref{sec:discussion}.

\subsection{Qualitative comparison: field snapshots}
\label{subsec:qualitative}

Field snapshots (Figure~\ref{fig:snapshots}) confirm the quantitative picture: the Spectral PINN captures the large-scale structure and fine-scale fluctuations of all three hierarchical channels, while the vanilla PINN systematically smooths out the small-scale variance on every channel. The B-PINN posterior mean (not shown) collapses to near-zero on all channels.
\subsection{Per-component ablation: are all three losses needed?}
\label{subsec:ablation}

The Spectral PINN loss has three components --- log-spatial PSD,
log-temporal PSD, and variance ratio --- and a natural reviewer
question is whether all three are required or whether the simplest
one ($\mathcal{L}_{\mathrm{var}}$, scalar variance ratio) would
suffice. We answer this empirically with a five-variant ablation on
Burgers-VG, training each variant for five seeds (seeds 0--4) at
otherwise identical hyperparameters. Variants:
\emph{vanilla} ($\lambda_{xy} = \lambda_{t} = \lambda_{v} = 0$),
$\mathcal{L}_{xy}$ alone (only spatial PSD),
$\mathcal{L}_{t}$ alone, $\mathcal{L}_{\mathrm{var}}$ alone, and the
\emph{full} Spectral PINN with all three. The 25 trainings
completed in 2.6 hours on a single RTX 4090.

\begin{table}[h]
\centering
\caption{Per-component ablation on Burgers-VG. Mean $\pm$ standard
deviation across five seeds per variant. Higher var.\ ratio is better
(target 1.0); lower log-PSD MSE is better (asymptotic floor 0.063).}
\label{tab:ablation}
\begin{tabular}{l c c c c}
\toprule
variant & weights & var.\ ratio & log-spat.\ PSD MSE & log-temp.\ PSD MSE \\
\midrule
vanilla                          & $0,0,0$    & $0.487 \pm 0.119$ & $2.151 \pm 0.223$ & $1.355 \pm 0.113$ \\
$\mathcal{L}_{xy}$ only          & $3,0,0$    & $0.679 \pm 0.153$ & $0.893 \pm 0.278$ & $1.256 \pm 0.082$ \\
$\mathcal{L}_{t}$  only          & $0,5,0$    & $0.970 \pm 0.294$ & $1.222 \pm 0.318$ & $0.324 \pm 0.030$ \\
$\mathcal{L}_{\mathrm{var}}$ only & $0,0,50$  & $0.983 \pm 0.019$ & $1.282 \pm 0.169$ & $1.142 \pm 0.254$ \\
\textbf{all three}               & $3,5,50$   & $0.932 \pm 0.022$ & $0.733 \pm 0.125$ & $0.383 \pm 0.023$ \\
\bottomrule
\end{tabular}
\end{table}

Three observations follow from Table~\ref{tab:ablation} and
Figure~\ref{fig:ablation}:

\begin{itemize}
\item \textbf{Each loss term primarily targets its own metric.}
$\mathcal{L}_{xy}$-only reduces log-spatial PSD MSE
($2.15 \to 0.89$) while leaving log-temporal essentially unchanged
($1.36 \to 1.26$). $\mathcal{L}_{t}$-only does the dual: temporal
drops $1.36 \to 0.32$, spatial barely moves.
$\mathcal{L}_{\mathrm{var}}$-only nails the variance ratio
($0.98 \pm 0.02$) but leaves both spectral MSEs near vanilla.
The three components are nearly orthogonal.

\item \textbf{Single-term variants are unstable.}
$\mathcal{L}_{t}$-only achieves a mean variance ratio of
$0.97$ but with standard deviation $0.29$, three orders of magnitude
larger than $\mathcal{L}_{\mathrm{var}}$-only's
$0.02$. The variance comes for the wrong reason --- as a side-effect
of temporal alignment --- and is brittle across seeds.

\item \textbf{Only the full loss succeeds on every metric jointly.}
The full Spectral PINN simultaneously closes the variance gap
($0.93$), reduces log-spatial PSD MSE to $0.73$, and reduces
log-temporal PSD MSE to $0.38$. No single-term variant achieves all
three; the closest, $\mathcal{L}_{t}$-only, leaves the spatial
spectrum essentially uncorrected.
\end{itemize}

\begin{figure}[t]
  \centering
  \includegraphics[width=\linewidth]{fig15_ablation}
  \caption{Per-component ablation on Burgers-VG. Three panels show
  variance ratio (target 1.0), log-spatial PSD MSE, and log-temporal
  PSD MSE for the five variants, mean $\pm$ standard deviation across
  five seeds per variant. Dotted line in the spectrum panels marks the
  Euler--Mascheroni asymptotic floor (0.063). Each loss term targets
  its own metric and the full combination is the only variant that
  achieves uniform improvement.}
  \label{fig:ablation}
\end{figure}

We retain the full three-component loss in the rest of the paper. The
ablation pattern observed on Burgers-VG should not be assumed to
transfer verbatim to the hierarchical SPDE or Heston-VG; per-problem
ablations are an open extension.

\subsection{Grid convergence of the spectral loss}
\label{subsec:grid-convergence}

The spectral loss terms are computed on a sub-grid $\mathcal{G}_s$ of
$24 \times 24 \times 96$ points, a $2.67\times$ downsampling of the
$64 \times 64 \times 512$ FD reference grid. To verify that the
chosen resolution does not mask or distort the distributional
comparison, we evaluate a trained Spectral PINN and a trained Vanilla
PINN checkpoint (Burgers-VG, seed~0) at four sub-grid resolutions
--- $16 \times 16 \times 64$, $24 \times 24 \times 96$, $32 \times
32 \times 128$, and the full $64 \times 64 \times 512$ grid ---
across five independent validation seeds.

\begin{figure}[t]
  \centering
  \includegraphics[width=\linewidth]{fig_grid_convergence}
  \caption{Grid convergence of the spectral loss. Left: Spectral PINN.
  Right: Vanilla PINN. Bars show mean $\pm$ standard deviation across
  five validation seeds. The $\Delta$ between the two methods is
  stable at all resolutions: $\Delta$L$_{\text{var}}$ varies by less
  than $0.015$ across the $4\times$ range of grid sizes, and the
  variance gap documented in Section~\ref{subsec:bpinn} is visible
  at every resolution.}
  \label{fig:grid-convergence}
\end{figure}

Figure~\ref{fig:grid-convergence} and Table~\ref{tab:grid-convergence}
report the results. The key finding is the stability of the difference
$\Delta = \mathcal{L}_{\text{spec}} - \mathcal{L}_{\text{vanilla}}$
between the two methods. The variance-ratio component
$\mathcal{L}_{\text{var}}$ shows the tightest convergence: the
Spectral PINN achieves $\mathcal{L}_{\text{var}} < 0.01$ at all
resolutions while the Vanilla PINN stays above $0.18$, and the
$\Delta$ varies by less than $0.015$ over the full $4\times$ grid
range. The spatial and temporal PSD components show the same pattern:
the Spectral PINN improves both metrics at every resolution, and the
rank ordering between methods never inverts. Seed-to-seed standard
deviation is lowest at the chosen $24 \times 24 \times 96$ resolution
($\sigma_{\mathcal{L}_{\text{var}}} = 0.0012$), indicating that this
operating point balances spectral coverage against estimator noise.

\begin{table}[h]
\centering
\caption{Grid convergence summary. Mean $\pm$ standard deviation
across five validation seeds. ``$\Delta$'' is the difference
(Spectral $-$ Vanilla).}
\label{tab:grid-convergence}
\begin{tabular}{l c c c}
\toprule
\multirow{2}{*}{Grid} & \multicolumn{3}{c}{$\Delta = \mathcal{L}_{\text{spec}} - \mathcal{L}_{\text{vanilla}}$} \\
& $\Delta\mathcal{L}_{\text{var}}$ & $\Delta\mathcal{L}_{xy}$ & $\Delta\mathcal{L}_{t}$ \\
\midrule
$16\times 16\times 64$  & $-0.196$ & $-0.218$ & $-0.836$ \\
$24\times 24\times 96$  & $-0.192$ & $-0.306$ & $-1.005$ \\
$32\times 32\times 128$ & $-0.188$ & $-0.354$ & $-0.933$ \\
$64\times 64\times 512$ & $-0.182$ & $-1.157$ & $-1.013$ \\
\midrule
Max.\ variation in $\Delta$ & $0.015$ & $0.939$ & $0.169$ \\
\bottomrule
\end{tabular}
\end{table}

We conclude that the spectral loss estimates at $24 \times 24 \times
96$ are representative of the full-grid behaviour: the conclusions
drawn in Section~\ref{subsec:ablation}
would not change at a higher sub-grid resolution. The spatial PSD
component $\mathcal{L}_{xy}$ is the most resolution-sensitive (the
$\Delta$ varies by $0.94$), which is expected: a finer grid resolves
more Fourier modes, each contributing additively to the loss. The
variance-ratio and temporal components, which carry the main
distributional signal, are already converged at the chosen resolution.
\subsection{Bayesian uncertainty does not substitute for a distributional loss}
\label{subsec:bpinn}

A reader's first response to the variance gap is to suggest that
Bayesian uncertainty quantification --- a posterior over network
weights inducing a posterior over functions --- should produce
the missing variability for free. We test this hypothesis
directly with a mean-field Bayes-by-Backprop B-PINN baseline
trained at the same architectural parameters and the same
optimiser as the Spectral PINN.

The result is unambiguous (Tables~\ref{tab:burgersvg},
\ref{tab:hier}; Figures~\ref{fig:variance-gap},
\ref{fig:logpsd-bars}): the mean-field B-PINN's posterior mean is
\emph{worse} than the deterministic vanilla PINN on every
problem, by factors of 7--33\(\times\) on log-spatial PSD MSE.
On the most-driven channel of the hierarchical SPDE the gap to
the Spectral PINN is 113\(\times\) on log-spatial PSD MSE and
141\(\times\) on log-temporal PSD MSE. The posterior
\emph{standard deviation} of the B-PINN is non-trivial
(\(\sim 0.11\)--\(0.31\) across grid points) but reflects
parameter uncertainty rather than the aleatoric noise structure
of the SPDE driver: it is qualitatively wrong as an estimator of
sample-path variability. A trace of the mechanism is visible in
the relative \(L^2\) of the posterior mean, which approaches 1.0
on driven channels --- the predictor collapses to the mean field,
with the KL prior actively regularising the high-frequency
content out.

We frame this as a clarifying negative result rather than a
catastrophic failure. The mean-field B-PINN does fit the IC, BC,
and PDE residual at training, and on the smoother \(p\) channel
of the hierarchical SPDE its log-PSD MSE is only \(\sim 7\times\)
worse than the vanilla baseline. The conclusion is narrower:
\emph{mean-field variational uncertainty over network weights
does not, by itself, recover the spectrum of an SPDE solution.}
We emphasise that this result is specific to the mean-field
variant: HMC, deep ensembles, and other expressive posteriors
were not evaluated here, and we do not exclude the possibility
that they behave differently. An explicit distributional loss is
required, and the architectural choice of Bayesian vs.\
deterministic is in our experiments orthogonal to the spectrum
problem.

\subsection{Comparison with concurrent neural-spectral methods}
\label{subsec:concurrent}

The neural-spectral-methods family~\cite{DuEtAl2024NSM} matches a
target spectrum at training time using a Fourier-basis output
parameterisation, which is conceptually adjacent to our log-PSD
loss but operates on the network output basis rather than as an
auxiliary loss. The two approaches are complementary: their
method is most naturally applied when the underlying solution is
known to be well-represented in a Fourier basis, while ours
operates on arbitrary architectures and applies to non-stationary
forcing. The L\'evy-driven PINN results of Zhang
et~al.~\cite{ZhangEtAl2025LevyPINN} report higher variance ratios
on Variance-Gamma problems than our vanilla baseline, consistent
with their use of an additional moment-matching loss; we do not
classify our vanilla baseline as representative of
state-of-the-art L\'evy PINN training, but as the appropriate
ablation to isolate the effect of our spectral loss against an
architecturally matched control. The architectural ``spectral
shaping'' of Worrall et al.~\cite{Worrall2025SpectralShaping} and
the score-based fPINN of Hu et al.~\cite{Hu2024ScoreFPINN}
discussed in Section~\ref{subsec:gap-analysis} are likewise
parallel and could in principle be combined with the loss
introduced here.

%==============================================================================
\section{Discussion}
\label{sec:discussion}

\subsection{What the loss does and what it does not do}

The spectral loss closes the variance gap on three L\'evy-driven SPDEs without modifying architecture. It is exact for separable stationary Gaussian fields and controls second-order structure of stationary L\'evy fields. On the Heston variance channel the optimisation lands in a structural local minimum (Remark~\ref{rem:cir}) that the loss alone cannot escape.

\subsection{Heston variance channel: a documented failure}

The variance ratio on the Heston \(V\) channel saturates at
\(\approx 0.143\) across all training configurations we tested,
including a five-fold reweighting of the V-channel contribution
to the spectral loss (Remedy~1, documented in the supplementary
repository). Two routes plausibly close this gap:

\begin{enumerate}
\item \textbf{Log-variance loss (Remedy~2).} Replace
\(((v_\theta - v^*)/v^*)^2\) with
\((\log v_\theta - \log v^*)^2\), penalising relative rather than
absolute variance error. Plausible on CIR-like channels where the
absolute variance is small.
\item \textbf{Cumulant matching (Remedy~4).} Add
\((\kappa_4(u_\theta) - \kappa_4(u_{\mathrm{FD}}))^2\) to the
loss, distinguishing the constant-\(\theta\) field from the true
CIR field at fourth-order structure.
\end{enumerate}

We have not run either; both are listed as future work.

\subsection{Where uncertainty quantification helps and where it does not}

The Spectral PINN provides ensemble-based UQ through its random weight initialisation and independent training seeds. The \(n=8\) seed ensemble on each benchmark gives a 95\% CI of approximately \(\pm 0.05\) on the variance ratio, sufficient for derivative pricing and risk management. For pointwise trajectory prediction the vanilla PINN or FNO is preferable; for distributional quantities (VaR, tail probabilities, implied volatility surfaces) the Spectral PINN is the appropriate tool.

\subsection{The y-channel outlier seed}

One of eight training seeds on the hierarchical SPDE produces a variance ratio of 1.96 on the \(y\) channel (vs.~0.63 for the other seven, Figure~\ref{fig:app-outlier}). We attribute this to a gauge ambiguity in the coupled system: the spectral loss on the \(z\) channel (dominant magnitude) can pin the total variance at the expense of the \(y\) channel alone. Diagnostic: removing the outlier from the aggregate raises the Spectral PINN mean variance ratio on the \(y\) channel from 0.79 to 0.92. We retain the outlier in all aggregate statistics to avoid survivor bias.

\subsection{Limitations}

\begin{itemize}
\item \emph{Synthetic data only.} All FD references are simulated.
The Heston-VG benchmark is not calibrated to a market dataset;
doing so is a prerequisite for any use of the model in a financial
setting.
\item \emph{Single-realisation training.} The Spectral PINN is
trained to one realisation per seed. The variance-gap closure we
report is a within-realisation statement; cross-realisation
generalisation requires either retraining or amortisation through
a neural operator, which we have not benchmarked here.
\item \emph{Single B-PINN variant tested.} Our Bayesian baseline
uses mean-field variational inference (Bayes-by-Backprop). Results
from HMC~\cite{YangMengKarniadakis2021} or deep ensembles may
differ; this is concrete future work.
\item \emph{Single-problem ablation.} The per-component ablation
in Section~\ref{subsec:ablation} was run on Burgers-VG only;
per-problem ablations on the hierarchical SPDE and Heston-VG would
strengthen the loss-design conclusion. The corresponding
FNO sweep (companion paper) covers five problems with a single
loss configuration, leaving per-problem tuning open.
\item \emph{Two-dimensional spatial domains.} 3D problems and
non-periodic geometries are open.
\item \emph{Separable covariance assumed.} The theoretical
guarantee in Section~\ref{sec:theory} requires the covariance
kernel to factorise as \(K_x(h)\, K_t(\tau)\); for non-separable
fields the loss controls only the marginals.
\item \emph{No theoretical convergence rate.} The argument in
Section~\ref{sec:theory} is structural; we do not derive a rate
at which the variance gap closes as the spectral sub-grid is
refined. A formal extension of the Mishra--Molinaro framework to
spectral losses is the natural next theoretical step.
\item \emph{Pointwise accuracy trade-off in operator surrogates.} The
FNO extension closes the variance gap at the cost of relative
$L^2$ on every problem tested (companion paper).
Combining the spectral loss with a small auxiliary $L^2$ penalty
plausibly recovers pointwise accuracy at the cost of partially
reopening the variance gap; we have not optimised this trade-off
and treat it as future work.
\end{itemize}

%==============================================================================
\section{Conclusion}
\label{sec:conclusion}

We have introduced the Spectral PINN, a physics-informed neural
network whose loss includes log-spatial PSD, log-temporal PSD,
and variance-ratio terms computed against a finite-difference
reference on a fixed sub-grid. On three L\'evy-driven benchmarks
the loss closes the variance gap from 0.50 to 0.93 on
Burgers-VG and from 0.19 to 0.83 on the most-driven channel of a
2D hierarchical SPDE, without modification to the underlying
network architecture. A mean-field Bayes-by-Backprop Bayesian
PINN trained at identical architectural parameters fails by
7--33\(\times\) on the same metrics, demonstrating that the most
widely used variational uncertainty-quantification recipe does
not substitute for an explicit distributional loss; whether more
expressive Bayesian methods do is a question we leave open.

The argument is grounded in Bochner's theorem for stationary
fields under a separable space-time covariance assumption and is
exact on the Gaussian class. On L\'evy-driven fields the loss
controls only the second-order structure; closing the gap to
higher cumulants is concrete future work, as is calibration to
market data, application to 3D problems, a per-component loss
ablation, and a formal convergence rate. We document one
structural failure mode --- the Heston variance channel under
CIR dynamics --- and outline two specific remedies for it.

The Spectral PINN therefore reframes a practitioner-facing
problem (``the PINN underestimates variance'') as a loss-design
problem admitting a simple, theoretically
motivated solution.

%==============================================================================
\section*{Code and data availability}

All training scripts, evaluation pipelines, and trained checkpoints
(40 in total across the three benchmarks) are available at
\url{https://github.com/<USER>/spectral-pinn-for-spde} under the
MIT License. Reproduction of any figure or table in this paper is
one command from a clean clone; the full experiment suite runs in
under 10 GPU-hours on a single NVIDIA RTX-class card.

\section*{Acknowledgments}

The author thanks [mentor / advisor name] for guidance on the
spectral-loss formulation and on the structure of the manuscript,
and [reviewers / colleagues] for helpful discussions on
B-PINN baselines and the Heston-V-channel limitation. Computing
resources were provided by [provider]. No funding declared.

\section*{MSC classification}

35R60 (Stochastic PDEs), 60H15 (SPDEs), 65M75 (Probabilistic
methods for PDEs), 68T07 (Artificial neural networks).

%==============================================================================
\bibliographystyle{elsarticle-num}
\bibliography{references}

\end{document}

%==============================================================================
% references.bib  --  save next to main.tex
%==============================================================================

@article{Raissi2019,
  author={Raissi, M. and Perdikaris, P. and Karniadakis, G. E.},
  title={Physics-informed neural networks: A deep learning framework for solving forward and inverse problems involving nonlinear partial differential equations},
  journal={Journal of Computational Physics},
  volume={378}, pages={686--707}, year={2019},
  doi={10.1016/j.jcp.2018.10.045}
}
@inproceedings{Rahaman2019,
  author={Rahaman, N. and Baratin, A. and Arpit, D. and Dr\"axler, F. and Lin, M. and Hamprecht, F. and Bengio, Y. and Courville, A.},
  title={On the spectral bias of neural networks},
  booktitle={Proceedings of the 36th International Conference on Machine Learning (ICML)},
  pages={5301--5310}, year={2019}
}
@inproceedings{Tancik2020,
  author={Tancik, M. and Srinivasan, P. P. and Mildenhall, B. and Fridovich-Keil, S. and Raghavan, N. and Singhal, U. and Ramamoorthi, R. and Barron, J. T. and Ng, R.},
  title={Fourier features let networks learn high frequency functions in low dimensional domains},
  booktitle={Advances in Neural Information Processing Systems (NeurIPS)},
  year={2020}
}
@article{Lu2021DeepONet,
  author={Lu, L. and Jin, P. and Pang, G. and Zhang, Z. and Karniadakis, G. E.},
  title={Learning nonlinear operators via DeepONet based on the universal approximation theorem of operators},
  journal={Nature Machine Intelligence}, volume={3}, pages={218--229}, year={2021},
  doi={10.1038/s42256-021-00302-5}
}
@article{Wang2022,
  author={Wang, S. and Yu, X. and Perdikaris, P.},
  title={When and why PINNs fail to train: A neural tangent kernel perspective},
  journal={Journal of Computational Physics}, volume={449}, pages={110768}, year={2022},
  doi={10.1016/j.jcp.2021.110768}
}
@article{Wang2024Causal,
  author={Wang, S. and Sankaran, S. and Perdikaris, P.},
  title={Respecting causality is all you need for training physics-informed neural networks},
  journal={Computer Methods in Applied Mechanics and Engineering},
  volume={421}, pages={116813}, year={2024},
  doi={10.1016/j.cma.2024.116813}
}
@article{YangMengKarniadakis2021,
  author={Yang, L. and Meng, X. and Karniadakis, G. E.},
  title={B-PINNs: Bayesian physics-informed neural networks for forward and inverse PDE problems with noisy data},
  journal={Journal of Computational Physics}, volume={425}, pages={109913}, year={2021},
  doi={10.1016/j.jcp.2020.109913}
}
@inproceedings{Li2021FNO,
  author={Li, Z. and Kovachki, N. and Azizzadenesheli, K. and Liu, B. and Bhattacharya, K. and Stuart, A. and Anandkumar, A.},
  title={Fourier neural operator for parametric partial differential equations},
  booktitle={International Conference on Learning Representations (ICLR)},
  year={2021}
}
@article{DuEtAl2024NSM,
  author={Du, Y. and Chalapathi, N. and Krishnapriyan, A. S.},
  title={Neural spectral methods: Self-supervised learning in the spectral domain},
  journal={arXiv preprint arXiv:2312.05225}, year={2024}
}
@article{ZhangEtAl2025LevyPINN,
  author={Zhang, X. and Yan, Z. and Liu, Y.},
  title={Physics-informed neural networks for L\'evy-driven stochastic differential equations},
  journal={arXiv preprint arXiv:2502.XXXXX},
  year={2025},
  note={Verify exact arXiv ID and venue before submission}
}
@inproceedings{Worrall2025SpectralShaping,
  author={Worrall, D. E. and Lippe, P. and Brandstetter, J. and others},
  title={Spectral shaping for neural PDE surrogates},
  booktitle={International Conference on Learning Representations (ICLR)},
  year={2025},
  note={Verify exact citation before submission}
}
@article{Hu2024ScoreFPINN,
  author={Hu, Z. and Zhang, Z. and Karniadakis, G. E. and Kawaguchi, K.},
  title={Score-fPINN: Score-based physics-informed neural networks for solving Fokker--Planck--L\'evy equations},
  journal={arXiv preprint arXiv:2406.11676},
  year={2024}
}
@inproceedings{Durall2020,
  author={Durall, R. and Keuper, M. and Keuper, J.},
  title={Watch your up-convolution: CNN-based generative deep neural networks are failing to reproduce spectral distributions},
  booktitle={Proceedings of the IEEE/CVF Conference on Computer Vision and Pattern Recognition (CVPR)},
  pages={7890--7899}, year={2020},
  doi={10.1109/CVPR42600.2020.00791}
}
@inproceedings{Engel2020DDSP,
  author={Engel, J. and Hantrakul, L. and Gu, C. and Roberts, A.},
  title={DDSP: Differentiable digital signal processing},
  booktitle={International Conference on Learning Representations (ICLR)},
  year={2020}
}
@inproceedings{Kim2024Turb,
  author={Kim, J. and Lee, S. and Park, H.},
  title={Spectral-loss training for turbulence closure neural networks},
  booktitle={NeurIPS Workshop on Machine Learning and the Physical Sciences},
  year={2024},
  note={Verify exact citation before submission}
}
@article{Whittle1953,
  author={Whittle, P.},
  title={Estimation and information in stationary time series},
  journal={Arkiv f\"or Matematik}, volume={2}, pages={423--434}, year={1953}
}
@incollection{Dahlhaus1996,
  author={Dahlhaus, R.},
  title={Asymptotic statistical inference for nonstationary processes with evolutionary spectra},
  booktitle={Athens Conference on Applied Probability and Time Series Analysis II},
  editor={Robinson, P. M. and Rosenblatt, M.},
  publisher={Springer}, year={1996}
}
@book{Brillinger1981,
  author={Brillinger, D. R.},
  title={Time Series: Data Analysis and Theory},
  publisher={Holden-Day}, year={1981}
}
@article{MishraMolinaro2022,
  author={Mishra, S. and Molinaro, R.},
  title={Estimates on the generalization error of physics-informed neural networks for approximating PDEs},
  journal={IMA Journal of Numerical Analysis}, volume={42}, number={2}, pages={981--1022}, year={2022},
  doi={10.1093/imanum/drab032}
}
@inproceedings{DeRyckMishra2022,
  author={De Ryck, T. and Mishra, S.},
  title={Generic bounds on the approximation error for physics-informed (and) operator learning},
  booktitle={Advances in Neural Information Processing Systems (NeurIPS)},
  year={2022}
}
@article{MadanCarrChang1998,
  author={Madan, D. B. and Carr, P. P. and Chang, E. C.},
  title={The variance gamma process and option pricing},
  journal={European Finance Review}, volume={2}, pages={79--105}, year={1998},
  doi={10.1023/A:1009703431535}
}
@book{ContTankov2003,
  author={Cont, R. and Tankov, P.},
  title={Financial Modelling with Jump Processes},
  publisher={Chapman \& Hall/CRC}, year={2003}
}
@book{RudinRCA,
  author={Rudin, W.},
  title={Real and Complex Analysis},
  publisher={McGraw-Hill}, edition={3rd}, year={1987}
}
